% BHU PYQ -- Manually transcribed & error-corrected
% Paper: MTB-201/MTB-202 : Statics and Dynamics (Old/New Course)
% Exam: B.Sc. (Hons.) Semester II Examination, 2013-14

\documentclass[11pt,a4paper]{article}
\usepackage[a4paper,top=2.5cm,bottom=2.5cm,left=2.8cm,right=2.8cm,headheight=14pt]{geometry}
\usepackage{amsmath,amssymb}
\usepackage[T1]{fontenc}\usepackage[utf8]{inputenc}
\usepackage{lmodern,microtype,fancyhdr,enumitem,hyperref,lastpage,parskip}

\hypersetup{colorlinks=true,urlcolor=black,linkcolor=black,
  pdftitle={MTB-202 Statics \& Dynamics -- BHU B.Sc. Sem II 2013-14}}
\pagestyle{fancy}\fancyhf{}
\renewcommand{\headrulewidth}{0.4pt}\renewcommand{\footrulewidth}{0.4pt}
\fancyhead[L]{\small   \textit{Banaras Hindu University}}
\fancyhead[R]{\small   \textit{MTB-201/202: Statics \& Dynamics}}
\fancyfoot[C]{\small Page  \thepage\ of \pageref{LastPage}}


\newlist{parts}{enumerate}{2}
\setlist[parts,1]{label=(\alph*),leftmargin=*,itemsep=5pt,topsep=3pt}
\setlist[parts,2]{label=(\roman*),leftmargin=*,itemsep=3pt,topsep=2pt}

\newcommand{\pts}[1]{\hfill   \textit{[#1]}}


\begin{document}

\begin{center}
  {\large\bfseries BANARAS HINDU UNIVERSITY}\\[2pt]
  {
\normalsize B.Sc.\ (Hons.) --- Semester II Examination, 2013-14}\\[6pt]
  \rule{0.8\linewidth}{0.5pt}\\[6pt]
  {\large\bfseries Paper No.\ MTB-201 (Old) / MTB-202 (New): Statics and Dynamics}\\[4pt]
  {
\normalsize Time: 3 Hours \hfill Full Marks: 70}\\[2pt]
  {\small Code:\ BS/Sem\,II/3}
\end{center}
\rule{\linewidth}{0.4pt}
\smallskip



\noindent   \textit{Instructions: Answer   \textbf{five} questions including Question~1 which is compulsory. Figures in right-hand margin indicate marks.}
\bigskip




\noindent   \textbf{Question 1.}   \textit{(Compulsory.)}

\begin{parts}
  \item Forces with components $(-150,0)$, $(30,40)$, $(0,-90)$, $(120,50)$ act at $(-6,1)$, $(-3,6)$, $(5,0)$, $(0,-5)$ respectively. Show they are in equilibrium.\pts{2}
  \item Three forces $P,Q,R$ act along sides $BC,CA,AB$ of $ riangle ABC$. If their resultant passes through the circumcentre, show $P\cos A+Q\cos B+R\cos C=0$.\pts{2}
  \item If the distance between two particles remains unchanged during a displacement, find the sum of works done by their mutual action and reaction.\pts{2}
  \item If the angular velocity of a moving point about a fixed origin is constant, show its transverse acceleration varies as its radial velocity.\pts{2}
  \item Radial and transverse velocity components of a particle are $\lambda r$ and $\mu \theta$. Find the path.\pts{2}
  \item If the path is a circle with the centre of force at the centre, find the law of force.\pts{2}
  \item State Kepler's second law. Show the transverse component of a planet's acceleration is zero.\pts{2}
  \item In SHM, show the average acceleration is obtained by multiplying its maximum by $\dfrac{2}{\pi}$.\pts{2}
  \item If a particle describes a circle with constant angular velocity, prove the foot of the perpendicular from the particle on any diameter executes SHM.\pts{3}
  \item A particle moves on a smooth curve under gravity with velocity proportional to the arc distance from the highest point. Prove the curve is a cycloid.\pts{3}
\end{parts}
\medskip\rule{\linewidth}{0.2pt}

\medskip


\noindent   \textbf{Question 2.}

\begin{parts}
  \item Show a system of coplanar forces can be reduced to a single force through a given point and a couple.\pts{6}
  \item Forces in the plane of $ riangle ABC$ are equivalent to a force at $A$ along the bisector of $\angle BAC$ and couple $G_1$. Moments about $B,C$ are $G_2,G_3$. Show $(b+c)G_1=bG_2+cG_3$.\pts{6}
\end{parts}
\medskip\rule{\linewidth}{0.2pt}

\medskip


\noindent   \textbf{Question 3.}

\begin{parts}
  \item State and prove the principle of virtual work for a system of coplanar forces.\pts{6}
  \item A string of length $a$ forms the shorter diagonal of a rhombus of four uniform rods each of length $b$ and weight $W$, hinged together, one rod horizontal. Prove the tension in the string is $\dfrac{2W(2b^2-a^2)}{b\sqrt{4b^2-a^2}}$.\pts{6}
\end{parts}
\medskip\rule{\linewidth}{0.2pt}

\medskip


\noindent   \textbf{Question 4.}

\begin{parts}
  \item Obtain the tangential and normal components of acceleration of a particle moving in a plane curve.\pts{6}
  \item If radial and transverse velocities are always proportional, and the same holds for accelerations, prove the velocity varies as some power of the radius vector.\pts{6}
\end{parts}
\medskip\rule{\linewidth}{0.2pt}

\medskip


\noindent   \textbf{Question 5.}

\begin{parts}
  \item A heavy bead slides on a smooth circular wire of radius $a$, projected from the lowest point with velocity just sufficient to reach the highest point. Show the radius through the bead turns through angle $2 an^{-1}\!\left(\sinh\!\left(t\sqrt{g/a}\right)\right)$ in time $t$.\pts{6}
  \item A particle slides down a smooth cycloid from rest at the cusp, axis vertical, vertex down. Show at the vertex the pressure on the curve equals twice the weight.\pts{6}
\end{parts}
\medskip\rule{\linewidth}{0.2pt}

\medskip


\noindent   \textbf{Question 6.}

\begin{parts}
  \item If the central orbit is $r^n=a^n\cos n \theta$ under force toward the pole, find the law of force.\pts{6}
  \item A particle moves under central acceleration $\mu/( \text{distance})^2$, projected with velocity $V$ at distance $R$. Show its path is a rectangular hyperbola if angle of projection is $\sin^{-1}\dfrac{\mu}{V\sqrt{V^2R^2-2\mu R}}$.\pts{6}
\end{parts}
\medskip\rule{\linewidth}{0.2pt}

\medskip


\noindent   \textbf{Question 7.}

\begin{parts}
  \item Discuss the motion of a particle in a frame moving in a straight line with acceleration relative to a fixed Newtonian frame.\pts{6}
  \item A particle performs SHM of period $T$ about centre $O$. It passes through $P$ (where $OP=b$) with velocity $v$ in the direction of $OP$. Show the time to return to $P$ is $\dfrac{T}{\pi} an^{-1}\!\dfrac{2\pi bv}{v^2T^2-4\pi^2b^2}$.\pts{6}
\end{parts}

\vfill\rule{\linewidth}{0.4pt}

\begin{center}\small
\href{https://bhu-exam.vercel.app/}{Click here to visit our website}\\[4pt]
\href{https://me-aryan.vercel.app/}{Click here to visit our contributor}\\[4pt]
\href{https://quashan-me.vercel.app/}{Click here to visit our contributor}
\end{center}
\end{document}
